% Options for packages loaded elsewhere
% Options for packages loaded elsewhere
\PassOptionsToPackage{unicode}{hyperref}
\PassOptionsToPackage{hyphens}{url}
\PassOptionsToPackage{dvipsnames,svgnames,x11names}{xcolor}
%
\documentclass[
  letterpaper,
  DIV=11,
  numbers=noendperiod]{scrartcl}
\usepackage{xcolor}
\usepackage{amsmath,amssymb}
\setcounter{secnumdepth}{-\maxdimen} % remove section numbering
\usepackage{iftex}
\ifPDFTeX
  \usepackage[T1]{fontenc}
  \usepackage[utf8]{inputenc}
  \usepackage{textcomp} % provide euro and other symbols
\else % if luatex or xetex
  \usepackage{unicode-math} % this also loads fontspec
  \defaultfontfeatures{Scale=MatchLowercase}
  \defaultfontfeatures[\rmfamily]{Ligatures=TeX,Scale=1}
\fi
\usepackage{lmodern}
\ifPDFTeX\else
  % xetex/luatex font selection
\fi
% Use upquote if available, for straight quotes in verbatim environments
\IfFileExists{upquote.sty}{\usepackage{upquote}}{}
\IfFileExists{microtype.sty}{% use microtype if available
  \usepackage[]{microtype}
  \UseMicrotypeSet[protrusion]{basicmath} % disable protrusion for tt fonts
}{}
\makeatletter
\@ifundefined{KOMAClassName}{% if non-KOMA class
  \IfFileExists{parskip.sty}{%
    \usepackage{parskip}
  }{% else
    \setlength{\parindent}{0pt}
    \setlength{\parskip}{6pt plus 2pt minus 1pt}}
}{% if KOMA class
  \KOMAoptions{parskip=half}}
\makeatother
% Make \paragraph and \subparagraph free-standing
\makeatletter
\ifx\paragraph\undefined\else
  \let\oldparagraph\paragraph
  \renewcommand{\paragraph}{
    \@ifstar
      \xxxParagraphStar
      \xxxParagraphNoStar
  }
  \newcommand{\xxxParagraphStar}[1]{\oldparagraph*{#1}\mbox{}}
  \newcommand{\xxxParagraphNoStar}[1]{\oldparagraph{#1}\mbox{}}
\fi
\ifx\subparagraph\undefined\else
  \let\oldsubparagraph\subparagraph
  \renewcommand{\subparagraph}{
    \@ifstar
      \xxxSubParagraphStar
      \xxxSubParagraphNoStar
  }
  \newcommand{\xxxSubParagraphStar}[1]{\oldsubparagraph*{#1}\mbox{}}
  \newcommand{\xxxSubParagraphNoStar}[1]{\oldsubparagraph{#1}\mbox{}}
\fi
\makeatother


\usepackage{longtable,booktabs,array}
\newcounter{none} % for unnumbered tables
\usepackage{calc} % for calculating minipage widths
% Correct order of tables after \paragraph or \subparagraph
\usepackage{etoolbox}
\makeatletter
\patchcmd\longtable{\par}{\if@noskipsec\mbox{}\fi\par}{}{}
\makeatother
% Allow footnotes in longtable head/foot
\IfFileExists{footnotehyper.sty}{\usepackage{footnotehyper}}{\usepackage{footnote}}
\makesavenoteenv{longtable}
\usepackage{graphicx}
\makeatletter
\newsavebox\pandoc@box
\newcommand*\pandocbounded[1]{% scales image to fit in text height/width
  \sbox\pandoc@box{#1}%
  \Gscale@div\@tempa{\textheight}{\dimexpr\ht\pandoc@box+\dp\pandoc@box\relax}%
  \Gscale@div\@tempb{\linewidth}{\wd\pandoc@box}%
  \ifdim\@tempb\p@<\@tempa\p@\let\@tempa\@tempb\fi% select the smaller of both
  \ifdim\@tempa\p@<\p@\scalebox{\@tempa}{\usebox\pandoc@box}%
  \else\usebox{\pandoc@box}%
  \fi%
}
% Set default figure placement to htbp
\def\fps@figure{htbp}
\makeatother





\setlength{\emergencystretch}{3em} % prevent overfull lines

\providecommand{\tightlist}{%
  \setlength{\itemsep}{0pt}\setlength{\parskip}{0pt}}



 


\KOMAoption{captions}{tableheading}
\makeatletter
\@ifpackageloaded{caption}{}{\usepackage{caption}}
\AtBeginDocument{%
\ifdefined\contentsname
  \renewcommand*\contentsname{Table of contents}
\else
  \newcommand\contentsname{Table of contents}
\fi
\ifdefined\listfigurename
  \renewcommand*\listfigurename{List of Figures}
\else
  \newcommand\listfigurename{List of Figures}
\fi
\ifdefined\listtablename
  \renewcommand*\listtablename{List of Tables}
\else
  \newcommand\listtablename{List of Tables}
\fi
\ifdefined\figurename
  \renewcommand*\figurename{Figure}
\else
  \newcommand\figurename{Figure}
\fi
\ifdefined\tablename
  \renewcommand*\tablename{Table}
\else
  \newcommand\tablename{Table}
\fi
}
\@ifpackageloaded{float}{}{\usepackage{float}}
\floatstyle{ruled}
\@ifundefined{c@chapter}{\newfloat{codelisting}{h}{lop}}{\newfloat{codelisting}{h}{lop}[chapter]}
\floatname{codelisting}{Listing}
\newcommand*\listoflistings{\listof{codelisting}{List of Listings}}
\makeatother
\makeatletter
\makeatother
\makeatletter
\@ifpackageloaded{caption}{}{\usepackage{caption}}
\@ifpackageloaded{subcaption}{}{\usepackage{subcaption}}
\makeatother
\usepackage{bookmark}
\IfFileExists{xurl.sty}{\usepackage{xurl}}{} % add URL line breaks if available
\urlstyle{same}
\hypersetup{
  pdftitle={Math 393 Syllabus},
  colorlinks=true,
  linkcolor={blue},
  filecolor={Maroon},
  citecolor={Blue},
  urlcolor={Blue},
  pdfcreator={LaTeX via pandoc}}


\title{Math 393 Syllabus}
\author{}
\date{}
\begin{document}
\maketitle


\section{Math 393: Abstract Algebra}\label{math-393-abstract-algebra}

\textbf{Fall 2026}\\
\textbf{Professor:} Geoffrey Buhl\\
\textbf{Email:}
\href{mailto:geoffrey.buhl@csuci.edu}{\nolinkurl{geoffrey.buhl@csuci.edu}}\\
\textbf{Phone:} 805-437-3122\\
\textbf{Office:} 2324 Marin Hall\\
\textbf{Office hours:} MW 9:00--10:00am or by appointment

\textbf{Course meetings:} Monday and Wednesday, 10:30--11:45am\\
\textbf{Course textbook:}
\href{https://judsonbooks.org/aata-files/aata-html/aata-toc.html}{Abstract
Algebra: Theory and Applications}, Thomas W. Judson\\
\textbf{Course website:} CI Learn

\subsection{Why I like this course}\label{why-i-like-this-course}

Discrete math is a fun departure from the calculus series, free from
concerns about continuity and the real numbers. We get our first chance
to explore abstract structures and a first introduction to proofs in
this course.

\subsection{Course Description}\label{course-description}

Prerequisite: Math 230 Topics include: Sets, algebraic systems, axioms,
definitions, propositions, and proofs. Combinatorics, graph theory,
moduli calculus. Coding, coding errors, and Hammang codes. Students are
expected to write mathematical proofs, and communicate mathematical
ideas clearly in written and oral form.

\subsection{How to contact me}\label{how-to-contact-me}

The best way to get in touch with me is in-person or by email. In
general, if you have any difficulties with the material, schedule, or
anything else, get in touch with me. I'm here to help, and it's much
easier to help if you are proactive. If my office hours don't work for
you, email her and we can work out a time to me in-person or via Zoom.

\subsection{Student Learning Outcomes}\label{student-learning-outcomes}

Upon successful completion of this course, students will be able to: -
Apply the principles of logic in a formal language, - Discuss the theory
and use of sets, - Compute the number of ways complicated tasks can be
performed, - Incorporate applied problems into graph theoretic framework
and using this interpretation to solve them, and - Express quantitative
ideas in oral and written form.

\subsection{Course Workload}\label{course-workload}

This is a 3 credit hour course, which means that students should be
prepared to spend a 9 hours a week on work related to this class outside
of class. Spending time throughout the week working on homework is
critical for a successful learning experience.

\subsection{Student Evaluation
Structure}\label{student-evaluation-structure}

{\def\LTcaptype{none} % do not increment counter
\begin{longtable}[]{@{}lr@{}}
\toprule\noalign{}
Component & Points \\
\midrule\noalign{}
\endhead
\bottomrule\noalign{}
\endlastfoot
Homework & 200 \\
Midterm & 100 \\
Proof Portfolio & 60 \\
Final Exam & 150 \\
\textbf{Total} & \textbf{650} \\
\end{longtable}
}

\subsection{Educational Opportunities}\label{educational-opportunities}

If you have special educational opportunities, like traveling for a
conference, going to Santa Rosa Island, participating in an academic
competition that would cause you to miss class, \emph{I strongly
encourage you to take part in these special opportunities.} Please let
me know in advance, and we can work together on a plan to address any
missed learning.

\subsection{Degree Planning}\label{degree-planning}

During Week 3 of classes of each semester, you are asked to update your
information in the degree planner. The degree planner is a useful tool
to help you stay on track toward achieving your goal of graduation.
Students can learn how to update the degree planner using information
provided here.

\subsection{In Case of Disruption}\label{in-case-of-disruption}

Wildfires and COVID are example of big disruptions that affect us all.
There are smaller ones, like the class I wasn't able to teach because my
second child was being born or when students have to help their family
members.~ Events big and small can affect our ability to be in class,
but the goal to make sure learning keeps happening.

So, yes disruptions big and small happen! A disruption means you, me, or
all of us cannot participate in `class as usual' for a reason we could
not predict at the beginning of the semester.~ After COVID-19, we are
well aware that disruption can happen fast. While we hope to avoid any
closures due to COVID related reasons, let's face it, disruption happens
every semester. You may get sick, I may get sick, or campus may close
due to a wildfire - this is California after all. Life happens!~ Our
goal as a learning community is to do our best to keep teaching and
learning with as little interruption as possible. This page explains
what you can expect from me and what I expect from you when facing
disruption - small or big!~

\textbf{If I am out}: If for whatever reason~I am unable to be on
campus, I will move lectures to a zoom format, which will be
communicated by the announcement in this course.~ I will record these
lecture and post them as well.

**If you are out*: Be proactive and communicate!~ I will always work
with you to figure out how to help you keep learning!~ I have old
lectures and notes that are recorded that I can share, if that is what
is needed.~ If you just need a little extra time to get something done,
let me know, I will always grant you a few more days on assignments.

**If campus closes*: If all of campus closes, which has happened at
least once over the last three or more academic years,~ we will move to
Zoom classes which I will also record and post in the pages for each
class.~ I will communicate changes to the course schedule or other
changes by using announcements in CILearn. More details can be found on
the ``Disruption Plan'' page for this course on CI Learn.

\subsection{Collaboration Policy}\label{collaboration-policy}

Working in groups on the homework assignment is strongly encouraged.
However each person should individually write up their solutions,
relying only on their notes about the problem. On each homework
assignment, please indicate which students you worked with. An important
part of mathematical culture and work is acknowledging the people that
helped you understand and solve a problem.

\subsection{Additional Resources}\label{additional-resources}

\textbf{Office hours} is a time when you can come to my office and ask
me questions about course material. Students who make use of office
house tend to do better in college courses. You are encouraged to make
early and regular use of \{\bf campus tutors\} in the LRC. For campus
tutoring locations, subjects and hours, go
\href{http://go.csuci.edu/tutoring}{this link}.

\subsection{Accomodations}\label{accomodations}

If you are a student with a disability requesting reasonable
accommodations in this course, please visit Disability Accommodations
and Support Services (DASS) located on the second floor of Arroyo Hall,
or call 805-437-3331. All requests for reasonable accommodations require
registration with DASS in advance of need. You can apply for DASS
services here. Faculty, students and DASS will work together regarding
classroom accommodations. You are encouraged to discuss approved
accommodations with your faculty.

\subsection{Civil Discourse Statement}\label{civil-discourse-statement}

All students, staff and faculty on our campus are expected to join in
making our campus a safe space for communication and civil discourse. If
you are experiencing discomfort related to the language you are hearing
or seeing on campus (in or out of classes), please talk with a trusted
faculty or staff member. Similarly, please consider whether the language
that you are using (in person or on Canvas) respects the rights of
others to engage in informed discourse and express a diversity of
opinions freely and in a civil manner (language from Academic Senate
Resolution SR 16-01, Commitment to Equity, Inclusion, and Civil
Discourse within our Diverse Campus Community).

\subsection{Basic Needs at CI}\label{basic-needs-at-ci}

If you face challenges securing food, housing, or other basic needs, you
are not alone, and CSUCI wants to help during this time of crisis. One
helpful resource is the community of staff available through the Basic
Needs Program (BNP) located on the first floor of Arroyo Hall, Room 114.
Students can call 805-437-2067, email basicneeds@csuci.edu, or drop in
during open hours and talk with a BNP student assistant or professional
staff member for resources, ideas, and strategies connected to basic
needs challenges. Students can complete a referral form to request
services for themselves or others by going to www.csuci.edu/basicneeds.
The BNP is known for the Dolphin Pantry located at Arroyo Hall 114, but
there are other resources available and staff who can help you work
through housing insecurity or displacement as well as financial
insecurity. Undergraduate students living in California are especially
encouraged to explore CalFresh (grocery money each month for eligible
students) as a resource. Domestic undergraduate students living in
California are likely to be eligible for CalFresh and BNP staff are
skilled with helping students navigate this process.

\subsection{Counseling and Psychological Services
(CAPS)}\label{counseling-and-psychological-services-caps}

CAPS is pleased to provide a wide range of services to assist students
in achieving their academic and personal goals. Services include
confidential short-term counseling, crisis intervention, psychiatric
consultation, and 24/7 phone counseling. CAPS is located in Bell Tower
East, 1867 and can be reached at 805-437-2088 (select option 1 on
voicemail for 24/7 phone counseling); you can also email us at
caps@csuci.edu or visit our website at https://www.csuci.edu/caps.

\subsection{Academic Honesty}\label{academic-honesty}

\emph{By enrolling at CSU Channel Islands, students are responsible for
upholding the University's policies and the Student Conduct Code.
Academic integrity and scholarship are values of the institution that
ensure respect for the academic reputation of the University, students,
faculty, and staff. Cheating, plagiarism, unauthorized collaboration
with another student, knowingly furnishing false information to the
University, buying, selling or stealing any material for an examination,
or substituting for another person may be considered violations of the
Student Conduct Code (located at
http://www.csuci.edu/campuslife/student-conduct/academic-dishonesty.htm).
If a student is found responsible for committing an act of academic
dishonesty in this course, the student may receive academic penalties
including a failing grade on an assignment or in the course, and a
disciplinary referral will be made and submitted to the Dean of Students
office. For additional information, please see the faculty Academic
Senate Policy on Academic Dishonesty, also in the CI Catalog. Please ask
about my expectations regarding academic dishonesty in this course if
they are unclear.}

Basically what this boils down to is: the work that you submit for
credit should be your own work and reflect your understanding of the
course material. Group work, collaborations, and the free exchange of
ideas are fundamental aspects of an academic environment. I encourage
you all to talk and work with others in the class, as it will enhance
your understanding of the material. You must write up your assignment
individually and explain solutions in your own words. It is also
important to cite resources you used in work, these may include other
students, books, websites, and instructors.

There is a distinction between collaboration and copying. Collaboration
is encouraged; copying solutions or assignments is not allowed. If I
find copying of solutions on an assignment it will result in a score of
0 on that assignment. Multiple instances of academic dishonesty will be
formally documented with the University.

\textbf{Academic Honesty and Generative AI:} I have begun to notice
students submitting work generated by AI. This is copying. If I suspect
AI has been used on an assignment, I will assign a score of 0, but you
can present your work in office hours or by appointment to receive
credit.

\subsection{Disclaimer on syllabus
changes}\label{disclaimer-on-syllabus-changes}

Information on this syllabus is subject to change. I will make such
changes known in class and on CI Learn.

\subsection{Semester Schedule}\label{semester-schedule}

comming soon




\end{document}
